\section{Acids, Bases, Buffers, Titration, and Solubility Products}

\subsection{Ctrl+F keywords: pH, pOH, weak acid, \(K_a\), \(pK_a\), buffer, Henderson, phosphate, titration, \(K_{sp}\), common ion, oploeselighed}
Use for acid strength, pH calculations, buffers, phosphate systems, titration equivalence, solubility, precipitation, and common-ion tasks.

\subsection{pH and acid-base constants}

\subsubsection{Water, pH, and pOH}
\[
K_w=[\ce{H3O+}][\ce{OH-}]=1.0\times10^{-14}, \qquad
\mathrm{pH}=-\log[\ce{H3O+}], \qquad
\mathrm{pH}+\mathrm{pOH}=14
\]
At \(298\,\mathrm K\), neutral water has pH \(7\). A decrease of one pH unit means ten times larger \([\ce{H3O+}]\).

\subsubsection{Strong acid and strong base pH}
\[
[\ce{H3O+}]\approx \nu_{\ce{H+}}c_{\mathrm{acid}}, \qquad
[\ce{OH-}]\approx \nu_{\ce{OH-}}c_{\mathrm{base}}
\]
Strong acids and bases are treated as fully dissociated. Include stoichiometric coefficients, for example \(\ce{Ba(OH)2}\) gives \(2\ce{OH-}\).

\subsubsection{Weak acid and weak base constants}
\[
K_a=\frac{[\ce{A-}][\ce{H3O+}]}{[\ce{HA}]}, \qquad
K_b=\frac{[\ce{HB+}][\ce{OH-}]}{[\ce{B}]}
\]
Larger \(K_a\) or smaller \(pK_a\) means stronger acid. Larger \(K_b\) or smaller \(pK_b\) means stronger base.

\subsubsection{Conjugate acid-base relation}
\[
K_aK_b=K_w, \qquad pK_a+pK_b=14
\]
Use this for conjugate acid/base pairs at \(298\,\mathrm K\).

\subsubsection{Weak acid and weak base approximations}
\[
[\ce{H3O+}]\approx\sqrt{K_ac_0}, \qquad
[\ce{OH-}]\approx\sqrt{K_bc_0}
\]
The approximation is valid when \(x/c_0\lesssim5\%\). Otherwise solve the quadratic.

\subsection{Exam recipe: pH of weak acid or weak base}
\textbf{Use when:} a weak acid/base concentration and \(K_a\), \(K_b\), \(pK_a\), or \(pK_b\) are given.

\textbf{Steps:}
\begin{enumerate}
    \item Write the equilibrium and initial concentration \(c_0\).
    \item Convert \(pK\) to \(K\) if needed: \(K=10^{-pK}\).
    \item Use \(x\approx\sqrt{Kc_0}\) for \([\ce{H3O+}]\) or \([\ce{OH-}]\).
    \item Convert to pH or pOH.
    \item Check \(x/c_0\); solve exactly if too large.
\end{enumerate}

\textbf{Common traps:} reporting pOH as pH for weak bases; forgetting \(pH=14-pOH\).

\subsection{Buffers and titration}

\subsubsection{Henderson-Hasselbalch buffer equation}
\[
\mathrm{pH}=pK_a+\log\left(\frac{[\ce{A-}]}{[\ce{HA}]}\right)
\]
Use mole ratios instead of concentrations when acid and base are in the same final volume.

\subsubsection{Buffer after adding strong acid or base}
\[
\ce{A- + H+ -> HA}, \qquad
\ce{HA + OH- -> A- + H2O}
\]
Do stoichiometric neutralization first, then apply Henderson-Hasselbalch to the remaining conjugate pair.

\subsubsection{Phosphoric acid phosphate system}
\[
\begin{aligned}
\ce{H3PO4 &<=> H+ + H2PO4-} && pK_{a1}\approx2.1,\\
\ce{H2PO4- &<=> H+ + HPO4^2-} && pK_{a2}\approx7.2,\\
\ce{HPO4^2- &<=> H+ + PO4^3-} && pK_{a3}\approx12.4.
\end{aligned}
\]
Choose the conjugate pair whose \(pK_a\) is closest to the target or observed pH.

\subsubsection{Titration equivalence point}
\[
n_{\mathrm{acid\ equivalents}}=n_{\mathrm{base\ equivalents}}
\]
At equivalence, stoichiometric acid-base equivalents have reacted. The pH depends on the salt left after neutralization.

\subsection{Exam recipe: buffer pH, phosphate buffer, Henderson}
\textbf{Use when:} both weak acid and conjugate base are present, or a buffer receives added strong acid/base.

\textbf{Steps:}
\begin{enumerate}
    \item Identify the conjugate pair and the relevant \(pK_a\).
    \item Convert all given concentrations/volumes to moles.
    \item If strong acid/base is added, neutralize stoichiometrically first.
    \item Insert remaining \(n_{\ce{A-}}/n_{\ce{HA}}\) into Henderson-Hasselbalch.
\end{enumerate}

\textbf{Fast checks:} when acid and base amounts are equal, \(\mathrm{pH}=pK_a\); buffer pH should be near \(pK_a\).

\subsection{Solubility products}

\subsubsection{Solubility product expression}
\[
\mathrm{M}_a\mathrm{X}_b(s)\rightleftharpoons a\,\mathrm{M}^{m+}+b\,\mathrm{X}^{x-}, \qquad
K_{sp}=[\mathrm{M}^{m+}]^a[\mathrm{X}^{x-}]^b
\]
The pure solid is omitted. Powers come from dissolution stoichiometry.

\subsubsection{Molar solubility patterns}
\[
\mathrm{MX}(s)\rightleftharpoons \mathrm{M}^{+}+\mathrm{X}^{-}: \qquad K_{sp}=s^2
\]
\[
\mathrm{MX_2}(s)\rightleftharpoons \mathrm{M}^{2+}+2\mathrm{X}^{-}: \qquad K_{sp}=s(2s)^2=4s^3
\]
\[
\mathrm{M(OH)_2}(s)\rightleftharpoons \mathrm{M}^{2+}+2\ce{OH-}: \qquad K_{sp}=s(2s)^2=4s^3
\]
Convert molar solubility \(s\) to ion concentrations using coefficients.

\subsubsection{Ion product and precipitation}
\[
Q_{sp}=[\mathrm{M}^{m+}]^a[\mathrm{X}^{x-}]^b
\]
\(Q_{sp}>K_{sp}\) means precipitation; \(Q_{sp}<K_{sp}\) means unsaturated; \(Q_{sp}=K_{sp}\) means saturated.

\subsubsection{Common-ion solubility}
\[
K_{sp}=[\mathrm{M}^{m+}][\mathrm{X}^{x-}], \qquad [\mathrm{X}^{x-}]_{\mathrm{initial}}\gg s \Rightarrow s\approx\frac{K_{sp}}{[\mathrm{X}^{x-}]_{\mathrm{initial}}}
\]
A common ion lowers solubility by shifting dissolution toward the solid.

\subsection{Exam recipe: \(K_{sp}\), molar solubility, common ion, hydroxide pH}
\textbf{Use when:} the problem asks for solubility, precipitation, saturated solution pH, or common-ion effect.

\textbf{Steps:}
\begin{enumerate}
    \item Write the dissolution reaction.
    \item Write \(K_{sp}\) with coefficient powers.
    \item Let \(s\) be molar solubility and express ion concentrations.
    \item Include any initial common-ion concentration before adding \(s\).
    \item For hydroxides, convert \([\ce{OH-}]\) to pOH and pH.
\end{enumerate}

\textbf{Common traps:} forgetting the square/cube from coefficients; using \(s\) instead of \(2s\) for hydroxide concentration.
